% 60-02(60쪽) — 삼차함수 $y=f(x)$ 의 그래프(왼쪽 위에서 내려와 극소 뒤 극대를 지나 오른쪽 아래로) · $x=a,\ b,\ c$ 의 안내 점선. 스캔 화질이 낮아 TikZ 로 다시 그림.
% 원문에 있는 것만: 곡선 · 세 안내 점선 · 라벨 a, O, b, c, y=f(x), x, y(눈금 수치·함수값은 원문에 없다).
\begin{tikzpicture}[x=0.52cm, y=0.28cm, line width=0.6pt]
  \draw[->, >=stealth, line width=0.7pt] (-1.50,0) -- (3.70,0) node[below=1pt] {$x$};
  \draw[->, >=stealth, line width=0.7pt] (0,-2.30) -- (0,7.60) node[left=1pt] {$y$};
  \draw[densely dotted, line width=0.4pt] (-0.85,3.78) -- (-0.85,0);
  \draw[densely dotted, line width=0.4pt] (1,3) -- (1,0);
  \draw[densely dotted, line width=0.4pt] (2,5) -- (2,0);
  \draw[domain=-1.15:3.25, samples=120, smooth] plot (\x, {-(\x)*(\x)*(\x)+3*(\x)*(\x)+1});
  \node[below=1pt, font=\small] at (-0.85,0) {$a$};
  \node[below right=-1pt and -2pt, font=\small] at (0,0) {$\pt{O}$};
  \node[below=1pt, font=\small] at (1,0) {$b$};
  \node[below=1pt, font=\small] at (2,0) {$c$};
  \node[anchor=south, font=\small] at (1.35,5.5) {$y=f(x)$};
\end{tikzpicture}
